\documentclass{article}
\usepackage[utf8]{inputenc}
\usepackage[T1]{fontenc}
\usepackage[a4paper,total={6in,8in}]{geometry}
\usepackage{amsmath,amssymb,amsfonts}
\usepackage{graphicx}
\usepackage{hyperref}
\usepackage{enumitem}
\usepackage{listings}
\usepackage{xcolor}
\DeclareMathOperator{\Var}{Var}
\DeclareMathOperator{\Cov}{Cov}
\providecommand{\tightlist}{\setlength{\itemsep}{0pt}\setlength{\parskip}{0pt}}
\lstset{basicstyle=\ttfamily\scriptsize,breaklines=true,columns=fullflexible,
  keepspaces=true,showstringspaces=false,frame=single}
\hypersetup{colorlinks=true,urlcolor=blue,linkcolor=blue}
\setlength{\emergencystretch}{2em}

\title{Math 126 Problem 1: Law of Large Numbers and Conditional Information}
\author{} % Put your name here.
\date{Spring 2027}
\begin{document}
\maketitle

\hypertarget{submission-instructions}{%
\subsection*{Submission Instructions}\label{submission-instructions}}

Please upload your work to Gradescope by \textbf{11:59pm on Monday, January 25, 2027 (Pacific time).}

\begin{itemize}
\item
  You are encouraged to discuss ideas and work with your classmates. However, you \emph{must acknowledge} your collaborators at the top of each solution on which you collaborated with others and you \emph{must write} your solutions independently.
\item
  Your solutions to theory questions must be written legibly, or typeset in LaTeX or markdown.
\item
  I recommend that you write your solutions in a Python notebook using Google Colab.
\item
  You should submit your solutions as a \emph{single PDF} via the assignment on Gradescope.
\end{itemize}

\textbf{Grading:} The point of the problem set is for \emph{you} to learn. To this end, I hope to disincentivize the use of LLMs by \emph{not} grading your work for correctness. Instead, you will grade your own work by comparing it to my solutions. This self-grade is due the Friday \emph{after} the problem set is due, also on Gradescope.

\noindent Original notebook: \href{https://colab.research.google.com/drive/1dVCKdYtpymvJlFcipDiUd813Liet0uWv?usp=sharing}{Google Colab}.

% Source notebook cell 2
\hypertarget{part-a-law-of-large-numbers}{%
\subsection*{Part A: Law of Large Numbers}\label{part-a-law-of-large-numbers}}

Let's consider the global distribution of heights: suppose it's normally distributed (a Gaussian) with mean \(\mu=5.63\) and standard deviation \(\sigma=.24\).

Generate \(n=1000\) samples from this distribution and plot a histogram.

Next, consider the distribution of \emph{average} heights in a sample of 100 people. Specifically, sample \(100\) people, compute their average height, and repeat \(n=1000\) times. Plot the distribution of these average heights on the same histogram as before.

How do the empirical expectations of the two distributions compare? How about their variances?

% Source notebook cell 3
\hypertarget{part-b-conditional-information}{%
\subsection*{Part B: Conditional Information}\label{part-b-conditional-information}}

Suppose I tell you that I have two cats. What's the probability that at least one of my cats is a boy when\ldots{}

\begin{itemize}
\item
  I don't provide any additional information.
\item
  I tell you that one of my cats is a girl (I haven't told you which).
\item
  I tell you that one of my cats is a girl and born on a Sunday (I still haven't told you which).
\end{itemize}

Why does the probability keep dropping?

\end{document}
